\documentclass[10pt,aspectratio=169]{beamer}

\usetheme{Madrid}
\usecolortheme{seahorse}
\setbeamertemplate{navigation symbols}{}

\usepackage[utf8]{inputenc}
\usepackage[T1]{fontenc}
\usepackage[french]{babel}
\usepackage{listings}
\usepackage{xcolor}
\usepackage{booktabs}
\usepackage{tikz}
\usetikzlibrary{arrows.meta, positioning}

\definecolor{codebg}{RGB}{245,245,245}
\definecolor{codegreen}{RGB}{40,120,60}
\definecolor{codeblue}{RGB}{30,80,160}

\lstset{
  basicstyle=\ttfamily\scriptsize,
  backgroundcolor=\color{codebg},
  keywordstyle=\color{codeblue}\bfseries,
  stringstyle=\color{codegreen},
  commentstyle=\color{gray},
  showstringspaces=false,
  breaklines=true,
  frame=single,
  framerule=0pt,
  language=Python,
}

\title{Serveur MCP d'outils d'analyse symbolique}
\subtitle{J3 --- IA symbolique : SAT/SMT, preuves par résolution, raisonnement OWL}
\author{Quentin Lauret}
\institute{EPITA}
\date{\today}

\begin{document}

% ---------------------------------------------------------------
\begin{frame}
  \titlepage
\end{frame}

% ---------------------------------------------------------------
\begin{frame}{Plan}
  \tableofcontents
\end{frame}

% ---------------------------------------------------------------
\section{Contexte et objectifs}

\begin{frame}{Contexte : pourquoi un serveur MCP symbolique ?}
  \begin{itemize}
    \item Les LLM raisonnent mal sur la logique formelle : hallucinations,
          erreurs de calcul, preuves invalides.
    \item Idée : \textbf{déléguer le raisonnement formel} à des outils
          symboliques exacts, exposés derrière une interface standard.
    \item Le \textbf{Model Context Protocol (MCP)} définit ce contrat :
          un client (LLM) découvre et appelle des outils via JSON-RPC.
  \end{itemize}
  \vfill
  \begin{block}{Objectif du projet}
    Implémenter un serveur JSON-RPC inspiré de MCP qui encapsule trois
    raisonneurs symboliques, avec gestion de sessions, exemples d'usage
    et évaluation de précision / robustesse.
  \end{block}
  \vfill
  \textit{Choix technique : Python pur, zéro dépendance obligatoire ---
  livrable pédagogique, reproductible et exécutable partout.}
\end{frame}

% ---------------------------------------------------------------
\section{Architecture}

\begin{frame}{Architecture générale}
  \begin{center}
  \begin{tikzpicture}[
      node distance=0.9cm and 1.4cm,
      box/.style={draw, rounded corners, align=center, minimum height=0.9cm,
                  minimum width=2.6cm, font=\small},
      tool/.style={box, fill=codebg},
      arr/.style={-{Stealth}, thick}
    ]
    \node[box, fill=blue!10] (client) {Client MCP\\(LLM outillé)};
    \node[box, fill=orange!15, right=of client] (server) {Serveur JSON-RPC\\\texttt{handle(request)}};
    \node[tool, right=of server] (sat) {SAT / SMT\\\scriptsize\texttt{symbolic.sat\_smt.solve}};
    \node[tool, below=0.35cm of sat] (proof) {Preuves\\\scriptsize\texttt{symbolic.proof.verify}};
    \node[tool, above=0.35cm of sat] (owl) {OWL-lite\\\scriptsize\texttt{symbolic.owl.reason}};
    \node[box, fill=green!10, below=of server] (session) {Sessions \& contexte\\\scriptsize historique des appels};

    \draw[arr] (client) -- node[above, font=\scriptsize]{JSON-RPC} (server);
    \draw[arr] (server) -- (sat);
    \draw[arr] (server) -- (proof);
    \draw[arr] (server) -- (owl);
    \draw[arr] (server) -- (session);
  \end{tikzpicture}
  \end{center}
  \vfill
  \begin{enumerate}
    \item Le client envoie une requête JSON-RPC (\texttt{initialize},
          \texttt{tools/list}, \texttt{tools/call}).
    \item Le serveur valide la méthode et la session, puis route vers l'outil.
    \item L'outil renvoie \texttt{structuredContent} (machine) +
          \texttt{content} (texte pour LLM).
    \item L'appel et son résultat sont mémorisés dans le contexte de session.
  \end{enumerate}
\end{frame}

% ---------------------------------------------------------------
\section{Les trois outils symboliques}

\begin{frame}[fragile]{Outil 1 --- Solveur SAT/SMT}
  \textbf{Deux familles d'entrées :}
  \begin{itemize}
    \item \textbf{SAT} : formule CNF en clauses d'entiers, résolue par
          \textbf{DPLL} (propagation unitaire + littéraux purs).
    \item \textbf{SMT borné} : contraintes arithmétiques sur entiers avec
          bornes finies, résolues par énumération.
  \end{itemize}
  \vfill
\begin{lstlisting}
# SAT : (P ou Q) et (~P ou Q) et (P ou ~Q)
{"clauses": [[1, 2], [-1, 2], [1, -2]], "symbols": ["P", "Q"]}
# -> sat, modele {P: true, Q: true}

# SMT borne
{"constraints": ["x + y == 7", "x >= 2", "y >= 2", "x < y"],
 "bounds": {"x": [0, 10], "y": [0, 10]}}
# -> sat, modele {x: 2, y: 5}
\end{lstlisting}
  \vfill
  \textbf{Sécurité :} les expressions SMT sont parsées avec le module
  \texttt{ast} --- aucun appel de fonction ni accès système autorisé.
\end{frame}

\begin{frame}[fragile]{Outil 2 --- Vérificateur de preuves}
  Vérifie des preuves par \textbf{résolution propositionnelle} : chaque
  étape doit dériver une clause par résolution de deux clauses précédentes.
  La clause vide \texttt{[]} prouve l'insatisfiabilité.
  \vfill
\begin{lstlisting}
# Premisses : (P ou Q), ~P, ~Q  -- refutation en 2 etapes
{"premises": [["P", "Q"], ["~P"], ["~Q"]],
 "proof": [
   {"id": "S1", "rule": "resolve", "from": ["P1", "P2"], "clause": ["Q"]},
   {"id": "S2", "rule": "resolve", "from": ["S1", "P3"], "clause": []}
 ],
 "conclusion": []}
# -> valid: true
\end{lstlisting}
  \vfill
  \begin{itemize}
    \item Notations acceptées : \texttt{"P"}, \texttt{"\textasciitilde P"},
          \texttt{"not P"}, $\neg$P, clauses en liste ou en chaîne.
    \item Chaque étape est \textbf{contrôlée formellement} : une résolution
          incorrecte invalide la preuve (pas de confiance aveugle).
  \end{itemize}
\end{frame}

\begin{frame}[fragile]{Outil 3 --- Raisonneur OWL-lite}
  Saturation de triplets \texttt{(sujet, prédicat, objet)} sur un fragment
  OWL/RDFS :
  \begin{columns}[T]
    \column{0.48\textwidth}
    \begin{itemize}
      \item transitivité de \texttt{subClassOf}, \texttt{subPropertyOf}
      \item propagation des types et assertions
      \item \texttt{equivalentClass}, \texttt{inverseOf}
      \item \texttt{domain}, \texttt{range}, \texttt{transitiveProperty}
      \item incohérences via \texttt{disjointWith}
    \end{itemize}
    \column{0.5\textwidth}
\begin{lstlisting}
["Dog", "subClassOf", "Mammal"]
["Mammal", "subClassOf", "Animal"]
["fido", "type", "Dog"]
["parentOf", "inverseOf", "childOf"]
["alice", "parentOf", "bob"]
# Inferences :
# fido type Animal
# bob childOf alice
\end{lstlisting}
  \end{columns}
  \vfill
  Point à saturation avec garde-fou (\texttt{max\_iterations}) ; le résultat
  distingue triplets de base et triplets \textbf{inférés}.
\end{frame}

% ---------------------------------------------------------------
\section{Interface MCP et orchestration LLM}

\begin{frame}[fragile]{Interface JSON-RPC / MCP}
  \begin{columns}[T]
    \column{0.48\textwidth}
    \textbf{Méthodes exposées :}
    \begin{itemize}
      \item \texttt{initialize}
      \item \texttt{tools/list} : schémas JSON des 3 outils
      \item \texttt{tools/call} : routage + validation
      \item \texttt{sessions/create}, \texttt{sessions/get}
    \end{itemize}
    \vspace{0.3cm}
    \textbf{Double sortie par outil :}
    \begin{itemize}
      \item \texttt{structuredContent} : exploitable par programme
      \item \texttt{content} : résumé textuel pour le LLM
    \end{itemize}
    \column{0.5\textwidth}
\begin{lstlisting}
{"jsonrpc": "2.0", "id": 4,
 "method": "tools/call",
 "params": {
   "session_id": "...",
   "name": "symbolic.sat_smt.solve",
   "arguments": {
     "clauses": [[1,2],[-1,2],[1,-2]],
     "symbols": ["P","Q"]}}}
\end{lstlisting}
  \end{columns}
  \vfill
  Chaque outil est documenté par un \texttt{inputSchema} JSON --- le client
  découvre dynamiquement les capacités du serveur.
\end{frame}

\begin{frame}{Traduction bidirectionnelle LLM $\leftrightarrow$ outils}
  \begin{itemize}
    \item \texttt{llm\_to\_tool\_call(message)} : un message en langage
          naturel est routé vers l'outil pertinent (détection d'intention
          simple par mots-clés : \emph{preuve}, \emph{ontologie},
          \emph{contrainte}\ldots).
    \item \texttt{tool\_result\_to\_llm\_answer(response)} : le résultat
          symbolique structuré est converti en réponse textuelle exploitable
          dans un dialogue.
    \item Le \textbf{contexte de session} conserve chaque appel : un
          orchestrateur peut réutiliser les résultats précédents, tracer
          les décisions et construire un raisonnement multi-étapes.
  \end{itemize}
  \vfill
  \begin{exampleblock}{Exemple}
    « Peux-tu vérifier cette ontologie OWL ? » \\
    $\Rightarrow$ appel \texttt{symbolic.owl.reason}
    $\Rightarrow$ « L'ontologie est cohérente, 6 triplets inférés\ldots »
  \end{exampleblock}
\end{frame}

% ---------------------------------------------------------------
\section{Évaluation}

\begin{frame}{Évaluation : précision et robustesse}
  \textbf{Précision} --- 8 benchmarks de référence, \alert{8/8 réussis (100\,\%)} :
  \vspace{0.2cm}
  \begin{center}
  \small
  \begin{tabular}{lll}
    \toprule
    \textbf{Benchmark} & \textbf{Outil} & \textbf{Attendu} \\
    \midrule
    sat\_cnf\_sat / sat\_cnf\_unsat & SAT & sat / unsat \\
    smt\_sat / smt\_unsat & SMT borné & modèle valide / unsat \\
    proof\_valid\_refutation & Preuves & preuve acceptée \\
    proof\_invalid\_step & Preuves & étape incorrecte rejetée \\
    owl\_type\_propagation & OWL & \texttt{fido type Animal} inféré \\
    owl\_disjoint\_inconsistency & OWL & incohérence détectée \\
    \bottomrule
  \end{tabular}
  \end{center}
  \vfill
  \textbf{Robustesse} --- entrées malveillantes ou invalides
  $\Rightarrow$ erreurs JSON-RPC propres (pas de crash) :
  outil inconnu, tentative d'injection
  \texttt{\_\_import\_\_('os').system(...)}, version JSON-RPC invalide.
\end{frame}

% ---------------------------------------------------------------
\section{Conclusion}

\begin{frame}{Limites et extensions}
  \begin{columns}[T]
    \column{0.48\textwidth}
    \textbf{Limites assumées}
    \begin{itemize}
      \item SMT borné sur entiers uniquement (énumération)
      \item Fragment OWL-lite, pas OWL DL complet
      \item Boucle de transport stdio non lancée dans le notebook
    \end{itemize}
    \column{0.48\textwidth}
    \textbf{Vers la production}
    \begin{itemize}
      \item Remplacer le mini-SMT par \textbf{Z3}
      \item Raisonneur OWL via \textbf{Owlready2} / RDFLib / OWL-RL
      \item Transport MCP complet (stdio ou HTTP)
    \end{itemize}
  \end{columns}
  \vfill
  \begin{block}{Bilan}
    Un serveur MCP fonctionnel qui montre comment donner à un LLM un accès
    fiable à du raisonnement symbolique exact : découverte d'outils,
    appels typés, sessions, et évaluation à 100\,\% de précision sur les
    benchmarks --- le tout en Python pur, sans dépendance.
  \end{block}
\end{frame}

\begin{frame}
  \begin{center}
    {\Huge Merci !}\\[0.8cm]
    {\large Questions ?}
  \end{center}
\end{frame}

\end{document}
